% ============================================================================
% PART 3: A MODEL ON YOUR OWN MACHINE
% ============================================================================

\sectionslide{Part 3. A Model on Your Own Machine}{Same task, same code, nothing leaves your laptop}

\begin{frame}{Why Run a Model Locally}

\small
\begin{columns}[T]
\column{0.48\textwidth}
\textbf{For}
\begin{itemize}
    \setlength{\itemsep}{4pt}
    \item \textbf{Confidential data stays on your machine}: interviews, survey answers, administrative records. Often the only option your data agreement allows.
    \item \textbf{Free per call}: a million sentences cost electricity and time, not money.
    \item \textbf{Exact replication}: you can archive the weights themselves, and cite their digest.
    \item \textbf{Temperature and seed work}: same input, same answer.
\end{itemize}

\column{0.48\textwidth}
\textbf{Against}
\begin{itemize}
    \setlength{\itemsep}{4pt}
    \item \textbf{Smaller, so less accurate}, on hard tasks above all (Halterman and Keith 2026: small open models and the Manifesto codebook).
    \item \textbf{Slower}, without a good GPU: seconds per sentence, not milliseconds.
    \item \textbf{Your hardware decides} which model you can run.
    \item \textbf{You maintain it}: versions, disk space, memory.
\end{itemize}
\end{columns}

\vspace{0.4cm}
\centering
{\small The middle way for large models on confidential data: the same open model on the UZH cluster (Part 4).}

\end{frame}

\begin{frame}[fragile]{Ollama in Five Commands}

\small
%%OLLAMACMDS%%

\end{frame}

\begin{frame}{Which Model for Your Machine}

\small
%%MODELTABLE%%

\end{frame}

\begin{frame}{Other Ways to Run a Model Locally}

\small
%%ALTERNATIVES%%

\end{frame}

\begin{frame}{Open Weights Is Not Open}

\small
%%OPEN%%

\end{frame}

\begin{frame}[fragile]{Same Code, Any Backend}

\small
Ollama, LM Studio, llama.cpp, vLLM and most commercial providers speak the same format (OpenAI's chat completions). Only three lines change.

\vspace{0.15cm}
{\scriptsize
\begin{verbatim}
from openai import OpenAI

client = OpenAI(base_url="http://localhost:11434/v1", api_key="not-needed")   # Ollama
# client = OpenAI(base_url="http://localhost:1234/v1",  api_key="not-needed")  # LM Studio
# client = OpenAI(base_url="http://localhost:8000/v1",  api_key="not-needed")  # vLLM on a GPU node
# client = OpenAI(base_url="https://api.mistral.ai/v1", api_key=MISTRAL_KEY)   # Mistral (EU)

r = client.chat.completions.create(
    model="%%PULLMODEL%%",
    messages=[{"role": "system", "content": CODEBOOK},
              {"role": "user", "content": "TARGET: ..."}],
    response_format={"type": "json_schema",
                     "json_schema": {"name": "annotation", "schema": SCHEMA, "strict": True}},
    temperature=0, seed=2026)
\end{verbatim}}

\vspace{0.1cm}
{\footnotesize In the notebook: \texttt{make\_openai\_annotator(base\_url, model)}, Part 5.}

\end{frame}

\begin{frame}{Block 3. Locally (11:05--11:30)}

\small
\textbf{I show (5 min).} Ollama on my laptop: \texttt{ollama ps}, the model in memory; Part 5 of the notebook; the digest in the bundle.

\vspace{0.3cm}

\textbf{You do (15 min)}
\begin{enumerate}
    \setlength{\itemsep}{3pt}
    \item \textbf{Laptop}: the notebook in VS Code, Cursor or Jupyter, with Ollama running. \textbf{Colab}: Runtime, Change runtime type, T4 GPU; then the ``Colab only'' cell.
    \item Part 5: one call, then the 200. Note the seconds per sentence.
    \item Part 6 again: two models now. Where do they disagree with each other?
    \item Faster than the others: pull a larger model, and compare.
\end{enumerate}

\vspace{0.3cm}

\textbf{We check (5 min).} The shared document: model and size, laptop or Colab, seconds per sentence, $\kappa$. Is the local model good enough for your project?

\end{frame}
